% Generated draft: requires mathematical and source review.
\documentclass[11pt]{article}
\usepackage[margin=25mm]{geometry}
\usepackage{fontspec}
\setmainfont{Latin Modern Roman}
\usepackage{amsmath,amssymb,mathtools,unicode-math}
\setmathfont{Latin Modern Math}
\usepackage{enumitem,needspace,xurl,hyperref}
\setlength{\parindent}{0pt}
\setlength{\parskip}{5pt}
\setlength{\emergencystretch}{3em}
\newcommand{\N}{\mathbb{N}}
\newcommand{\Z}{\mathbb{Z}}
\newcommand{\Q}{\mathbb{Q}}
\newcommand{\R}{\mathbb{R}}
\newcommand{\C}{\mathbb{C}}
\newcommand{\F}{\mathbb{F}}
\newcommand{\A}{\mathbb{A}}
\newcommand{\PP}{\mathbb{P}}
\newcommand{\E}{\mathbb{E}}
\newcommand{\eps}{\varepsilon}
\newcommand{\dd}{\,\mathrm{d}}
\begin{document}
\section*{UnsolvedMath ID: 2000018}
\textbf{Source problem code:} catalogue entry
\subsection{Definitions and mathematical statement}
Let $h>0$ denote thickness and let $\Omega_h\subset\R^3$ be a family of thin reference bodies. Rods shrink in two transverse directions, plates shrink in one direction about a planar mid-surface, and shells shrink about a curved mid-surface. For example, a plate family can have the form $\omega\times(-h/2,h/2)$ for a planar region $\omega$; this is illustrative, not a restriction supplied in the source.

A representative three-dimensional elastic energy is
\[
 E_h(y)=\int_{\Omega_h}W(x,\nabla y(x))\,dx-\mathcal L_h(y),
\]
where $y:\Omega_h\to\R^3$ is an admissible deformation, $W$ is stored energy, and $\mathcal L_h$ represents loads. After a specified rescaling to a fixed reference domain, one seeks an effective one-dimensional rod functional or two-dimensional plate or shell functional as $h\downarrow0$. The energy normalization, loads, material assumptions, and limiting variables depend on the regime and are not specified in the record.

\textbf{Problem.} Give rigorous derivations of models of rods, plates, and shells from three-dimensional elasticity as thickness tends to zero. A variational version, as discussed in the supplied progress summary, is to identify effective functionals $E_0$ and normalizations $c_h$ for which rescaled $c_hE_h$ $\Gamma$-converge in a specified topology: every convergent sequence obeys the lower bound
\[
 E_0(u)\leq\liminf_{h\downarrow0}c_hE_h(y_h),
\]
and every admissible limit $u$ has a recovery sequence with
\[
 \limsup_{h\downarrow0}c_hE_h(y_h)\leq E_0(u).
\]
Here convergence must be defined after identification of the varying domains and deformations with the chosen limiting variables. Compactness/coercivity assumptions are also needed to obtain conclusions about minimizing sequences. This is a description of a rigorous derivation framework, not a claim that the source fixes one scaling, topology, or particular model.

The supplied research summary calls plate and rod derivations solved in the literature, citing Friesecke--James--M\"uller (2002) and follow-ups, and says many shell regimes have been derived rigorously. It leaves open aspects of a unified full nonlinear hierarchy for arbitrary shell geometry. Its broad solved labels do not specify unrestricted material, loading, or boundary hypotheses and are not independently verified here.
\subsection{Short English statement}
Derive lower-dimensional rod, plate, and shell theories as limits of three-dimensional elasticity when thickness vanishes. The summary reports established variational derivations for rods and plates and many shell regimes, while leaving a fully unified nonlinear shell theory incomplete. The exact scalings and hypotheses are absent.
\subsection{Sources}
\begin{itemize}
\item Local UnsolvedMath registry, permanent ID 2000018, source code catalogue entry, ``Some Open Problems in Elasticity --- Dimension reduction''; status partially solved, classification PARTIAL-PROGRESS, updated 2026-08-06.
\item Ball, source list ``problems on Elasticity'' (2002), Problem 18, \url{https://people.maths.ox.ac.uk/ball/Articles%20in%20Conference%20Proceedings%20and%20Books/JMB%202002%20re%20Marsden%2060th.pdf}. The record reports PDF-text extraction and access on 2026-07-29. It says the source presents an open item and current status requires release review. Rights: NEEDS\_REVIEW; public author-hosted PDF.
\item The supplied research summary cites Friesecke--James--M\"uller (2002) and unspecified follow-ups for rigorous $\Gamma$-convergence derivations, and describes partial progress on general shells. No article titles, precise theorem numbers, or complete bibliography are supplied.
\end{itemize}
\end{document}
